% ============================================================
%  PEMBAHASAN SUPER DETAIL — LATIHAN SOAL & STUDI KASUS
%  BAB 8: PELUANG
%  Catatan: Soal Akhir Bab TIDAK dibahas di sini (sesuai permintaan).
% ============================================================
% ============================================================
%  PREAMBUL BERSAMA — BUKU PEMBAHASAN LATIHAN SOAL
%  Dipakai oleh MAE-2026-2027-Latihan1.tex ... Latihan9.tex
%  Kompilasi: pdfLaTeX (disarankan Overleaf / MiKTeX)
% ============================================================
\documentclass[11pt,a4paper]{report}

% ---------- PAKET ----------
\usepackage[utf8]{inputenc}
\usepackage[T1]{fontenc}
\usepackage[bahasa]{babel}
\usepackage{amsmath,amssymb,amsthm,mathtools}
\usepackage{geometry}
\usepackage{xcolor}
\usepackage{graphicx}
\usepackage{enumitem}
\usepackage{tikz}
\usetikzlibrary{calc,arrows.meta,angles,quotes,positioning,decorations.pathreplacing,patterns}
\usepackage{pgfplots}
\pgfplotsset{compat=1.18}
\usepackage[most]{tcolorbox}
\usepackage{fancyhdr}
\usepackage{titlesec}
\usepackage{array,booktabs}
\usepackage{multicol}
\usepackage{pifont}
\usepackage[colorlinks=true,linkcolor=biru,urlcolor=biru]{hyperref}

\geometry{a4paper,margin=2.5cm,headheight=15pt}

% ---------- WARNA ----------
\definecolor{biru}{HTML}{1F4E79}
\definecolor{birumuda}{HTML}{D6E4F0}
\definecolor{hijau}{HTML}{2E7D32}
\definecolor{hijaumuda}{HTML}{E3F2E1}
\definecolor{oranye}{HTML}{C75B12}
\definecolor{oranyemuda}{HTML}{FBE8DA}
\definecolor{abu}{HTML}{F2F2F2}
\definecolor{ungu}{HTML}{6A1B9A}
\definecolor{ungumuda}{HTML}{EDE2F3}
\definecolor{merah}{HTML}{C62828}
\definecolor{merahmuda}{HTML}{FCE4E4}
\definecolor{tosca}{HTML}{00838F}
\definecolor{toscamuda}{HTML}{E0F2F1}
\definecolor{emas}{HTML}{8D6E00}
\definecolor{emasmuda}{HTML}{FFF6D6}

% ---------- HEADER / FOOTER ----------
\newcommand{\babjudul}{Pembahasan Latihan}
\pagestyle{fancy}
\fancyhf{}
\renewcommand{\headrulewidth}{0.8pt}
\renewcommand{\footrulewidth}{0pt}
\fancyhead[L]{\small\textcolor{ungu}{\leftmark}}
\fancyhead[R]{\small\textcolor{ungu}{\babjudul}}
\fancyfoot[C]{\thepage}

% ---------- GAYA JUDUL BAB ----------
\titleformat{\chapter}[display]
  {\normalfont\huge\bfseries\color{ungu}}
  {\filright\large PEMBAHASAN}{8pt}{\Huge}
\titlespacing*{\chapter}{0pt}{0pt}{20pt}

% ---------- LINGKUNGAN TEOREMA ----------
\newtheoremstyle{gaya}{6pt}{6pt}{}{}{\bfseries\color{biru}}{.}{.5em}{}
\theoremstyle{gaya}
\newtheorem{definisi}{Definisi}[chapter]
\newtheorem{sifat}{Sifat}[chapter]

% ---------- KOTAK: SOAL (ungu) ----------
\newtcolorbox{soalbox}[1]{
  enhanced, breakable, colback=ungumuda, colframe=ungu, boxrule=0.8pt,
  arc=3pt, left=10pt, right=10pt, top=8pt, bottom=8pt,
  title={\textbf{Soal #1}}, fonttitle=\bfseries\color{white}, coltitle=white,
  attach boxed title to top left={xshift=10pt,yshift=-10pt},
  boxed title style={colback=ungu,arc=2pt}}

% ---------- KOTAK: KONSEP KUNCI (hijau) ----------
\newtcolorbox{ingat}[1]{
  enhanced, breakable, colback=hijaumuda, colframe=hijau, boxrule=0.6pt,
  arc=3pt, left=10pt, right=10pt, top=6pt, bottom=6pt,
  title={\textbf{Konsep Kunci: #1}}, fonttitle=\bfseries\color{white}, coltitle=white,
  attach boxed title to top left={xshift=10pt,yshift=-10pt},
  boxed title style={colback=hijau,arc=2pt}}

% ---------- KOTAK: JEBAKAN (merah) ----------
\newtcolorbox{jebakan}{
  enhanced, breakable, colback=merahmuda, colframe=merah, boxrule=0.6pt,
  arc=3pt, left=10pt, right=10pt, top=6pt, bottom=6pt,
  title={\textbf{\ding{55}\ Jebakan \& Kesalahan yang Sering Terjadi}},
  fonttitle=\bfseries\color{white}, coltitle=white,
  attach boxed title to top left={xshift=10pt,yshift=-10pt},
  boxed title style={colback=merah,arc=2pt}}

% ---------- KOTAK: VARIASI SOAL (oranye) ----------
\newtcolorbox{variasi}{
  enhanced, breakable, colback=oranyemuda, colframe=oranye, boxrule=0.6pt,
  arc=3pt, left=10pt, right=10pt, top=6pt, bottom=6pt,
  title={\textbf{\ding{72}\ Bentuk Modifikasi Soal Ini}},
  fonttitle=\bfseries\color{white}, coltitle=white,
  attach boxed title to top left={xshift=10pt,yshift=-10pt},
  boxed title style={colback=oranye,arc=2pt}}

% ---------- KOTAK: STUDI KASUS (biru) ----------
\newtcolorbox{studikasus}[1]{
  enhanced, breakable, colback=abu, colframe=biru, boxrule=0.8pt,
  arc=3pt, left=10pt, right=10pt, top=8pt, bottom=8pt,
  title={\textbf{Studi Kasus: #1}}, fonttitle=\bfseries\color{white}, coltitle=white,
  attach boxed title to top left={xshift=10pt,yshift=-10pt},
  boxed title style={colback=biru,arc=2pt}}

% ---------- KOTAK: TIPS (tosca) ----------
\newtcolorbox{tips}{
  enhanced, breakable, colback=toscamuda, colframe=tosca, boxrule=0.6pt,
  arc=3pt, left=10pt, right=10pt, top=6pt, bottom=6pt,
  borderline west={3pt}{0pt}{tosca}}

% ---------- PERINTAH BANTU ----------
\newcommand{\jawab}{\par\smallskip\noindent\textit{\textcolor{oranye}{Penyelesaian langkah demi langkah:}}\par\smallskip}
\newcommand{\langkah}[1]{\par\smallskip\noindent\textbf{\textcolor{biru}{Langkah #1.}}\ }
\newcommand{\jadi}{\par\smallskip\noindent$\Rightarrow$\ \textbf{Jadi, }}
% Judul tiap nomor soal
\newcommand{\soalno}[2]{%
  \par\bigskip
  \noindent{\large\bfseries\color{ungu}\rule[-2pt]{4pt}{14pt}\ Soal #1}%
  \ {\normalfont\itshape\color{gray}(#2)}\par\nopagebreak\smallskip}

% ---------- HALAMAN JUDUL OTOMATIS ----------
% #1 = judul topik (mis. "Barisan dan Deret"); #2 = nomor bab (angka)
\newcommand{\buathalamanjudul}[2]{%
  \begin{titlepage}
  \centering
  \vspace*{2cm}
  {\color{ungu}\rule{\textwidth}{2pt}}\\[0.6cm]
  {\Huge\bfseries\color{ungu} PEMBAHASAN\\[0.3cm] LATIHAN SOAL}\\[0.4cm]
  {\Large BAB #2 --- #1}\\[0.2cm]
  {\color{ungu}\rule{\textwidth}{2pt}}\\[1.2cm]

  {\large Matematika SMA --- Fase E (Kelas 10)}\\[0.2cm]
  {\large Kurikulum Merdeka}\\[2cm]

  \begin{tcolorbox}[enhanced,width=0.85\textwidth,colback=ungumuda,colframe=ungu,
    arc=4pt,boxrule=1pt]
  \centering\small
  Buku pembahasan ini mengupas \textbf{setiap soal pada kotak ungu (Latihan Soal)}
  di tiap subbab, \emph{selangkah demi selangkah seperti mengajari dari nol}:
  konsep yang dipakai, jebakan yang sering menjerat, cara soal bisa dimodifikasi,
  dan visualisasi. \emph{Soal Akhir Bab tidak dibahas di sini.}
  \end{tcolorbox}

  \vfill
  {\large Tahun Pelajaran 2026/2027}
  \end{titlepage}
  \tableofcontents
  \thispagestyle{empty}
  \clearpage}

% ============================================================
%  PENJAGA KOMPILASI
%  File ini adalah PREAMBUL BERSAMA (di-\input oleh Latihan1..9),
%  BUKAN untuk dikompilasi sendiri. Jika editor terlanjur mengompilasi
%  file ini secara langsung, tampilkan satu halaman info (bukan error).
%  Saat di-\input oleh berkas Latihan, blok ini dilewati.
% ============================================================
\ifnum\pdfstrcmp{\jobname}{MAE-2026-2027-Latihan-preamble}=0
  \begin{document}
  \thispagestyle{empty}
  \begin{center}
  \vspace*{3cm}
  {\Huge\bfseries\color{ungu} Berkas Preambel Bersama}\\[1cm]
  \begin{tcolorbox}[enhanced,width=0.85\textwidth,colback=ungumuda,
    colframe=ungu,arc=4pt,boxrule=1pt]
  \centering
  File ini \textbf{bukan untuk dikompilasi sendiri.} Ia berisi pengaturan
  bersama (paket, warna, kotak, perintah) yang dipanggil lewat
  \texttt{\textbackslash input} oleh berkas \textbf{MAE-2026-2027-Latihan1.tex}
  sampai \textbf{Latihan9.tex}.

  \medskip
  Silakan buka dan kompilasi salah satu berkas \textbf{Latihan1--Latihan9}
  tersebut untuk menghasilkan PDF pembahasan.
  \end{tcolorbox}
  \end{center}
  \end{document}
\fi

\renewcommand{\babjudul}{Pembahasan Latihan — Bab 8}

\begin{document}
\setcounter{chapter}{8}
\buathalamanjudul{Peluang}{8}

% ============================================================
\chapter*{Pembahasan Bab 8: Peluang}
\addcontentsline{toc}{chapter}{Pembahasan Bab 8: Peluang}
\markboth{PEMBAHASAN BAB 8}{}
% ============================================================

\begin{tips}
\textbf{Cara membaca buku ini.} Tiap soal: \textbf{(1)} soal ditulis ulang,
\textbf{(2)} konsep kunci, \textbf{(3)} langkah demi langkah, \textbf{(4)}
jebakan, \textbf{(5)} variasi. \emph{Aturan emas Bab 8:} selalu tuliskan
\textbf{ruang sampel} (atau hitung $n(S)$) lebih dulu, dan tanya ``apakah urutan
penting?'' serta ``apakah dengan/tanpa pengembalian?''.
\end{tips}

% ============================================================
\section*{Studi Kasus: Prediksi Cuaca dan Peluang Hujan}
\addcontentsline{toc}{section}{Studi Kasus: Peluang Hujan}
% ============================================================

\begin{studikasus}{Prediksi Cuaca dan Peluang Hujan}
Selama $100$ hari berkondisi ``berawan tebal'', hujan terjadi pada $70$ hari.
\textbf{Pertanyaan.} (a) Bagaimana menaksir ``peluang hujan'' dari data ini?
(b) Apakah peluang ini pasti benar untuk besok? (c) Apa bedanya dengan peluang
dadu yang bisa dihitung tanpa percobaan?
\end{studikasus}

\subsection*{Jawaban lengkap studi kasus}

\begin{ingat}{Peluang empiris vs teoretis}
\textbf{Empiris} (frekuensi relatif) $=\dfrac{\text{banyak kejadian}}{\text{banyak
percobaan}}$ --- diukur dari data. \textbf{Teoretis} $=\dfrac{n(A)}{n(S)}$ ---
dihitung dari asumsi hasil ``sama mungkin'', tanpa percobaan.
\end{ingat}

\textbf{(a) Menaksir peluang hujan.} Gunakan frekuensi relatif:
\[
P(\text{hujan})\approx\frac{70}{100}=0{,}7=70\%.
\]

\textbf{(b) Apakah pasti untuk besok?} \textbf{Tidak.} Ini \emph{taksiran}
berdasarkan pola masa lalu, bukan jaminan. ``$70\%$'' berarti: dari banyak hari
berkondisi serupa, sekitar $70\%$ berujung hujan --- besok tetap tak pasti.

\textbf{(c) Beda dengan dadu.} Peluang dadu bersifat \emph{teoretis}: kita tahu
$P(\text{mata }3)=\frac16$ tanpa melempar, karena keenam sisi sama mungkin.
Peluang cuaca bersifat \emph{empiris}: tak ada rumus ``sisi sama mungkin'',
sehingga harus \emph{diamati} dari data.
\jadi $P(\text{hujan})\approx0{,}7$ (empiris), tak pasti untuk satu hari, dan
berbeda asal-usulnya dari peluang dadu (teoretis).

\begin{tips}
\textbf{Hukum bilangan besar:} makin banyak percobaan, peluang empiris makin
mendekati teoretis. Inilah alasan kasino dan asuransi untung dalam jangka panjang
meski tiap kejadian tampak acak.
\end{tips}

% ############################################################
\section*{Subbab 8.1 — Konsep Dasar Peluang}
\addcontentsline{toc}{section}{Subbab 8.1 — Konsep Dasar}
% ############################################################

\begin{ingat}{Rumus \& aturan dasar}
$P(A)=\dfrac{n(A)}{n(S)}$, $0\le P(A)\le1$. Komplemen: $P(A^c)=1-P(A)$.
Saling lepas: $P(A\cup B)=P(A)+P(B)$. Tidak saling lepas:
$P(A\cup B)=P(A)+P(B)-P(A\cap B)$. Saling bebas: $P(A\cap B)=P(A)P(B)$.
\end{ingat}

% ---------------- SOAL 8.1.1 ----------------
\soalno{8.1.1}{ruang sampel dua koin}
\begin{soalbox}{}
Dua koin dilempar bersamaan. Tentukan peluang muncul tepat satu Angka.
\end{soalbox}

\jawab
\langkah{1} Tuliskan ruang sampel: $S=\{AA,\ AG,\ GA,\ GG\}$, jadi $n(S)=4$.
\langkah{2} Kejadian ``tepat satu Angka'': $\{AG,\ GA\}$, jadi $n(A)=2$.
\langkah{3} $P=\dfrac{n(A)}{n(S)}=\dfrac{2}{4}=\dfrac12$.
\jadi peluangnya $\dfrac12$.

\begin{jebakan}
\begin{itemize}[leftmargin=*,topsep=2pt]
  \item \textbf{Menganggap $AG$ dan $GA$ sama.} Koin pertama-A koin kedua-G
  \emph{berbeda} dari sebaliknya; keduanya dihitung. Kalau digabung, penyebutnya
  salah.
  \item \textbf{``Tepat satu'' vs ``minimal satu''.} Tepat satu = $2$ kasus;
  minimal satu Angka = $3$ kasus ($AA,AG,GA$).
\end{itemize}
\end{jebakan}

\begin{variasi}
Tepat dua Angka ($\frac14$); minimal satu Angka ($\frac34$, atau lewat komplemen
$1-P(GG)$).
\end{variasi}

% ---------------- SOAL 8.1.2 ----------------
\soalno{8.1.2}{peluang \& komplemen}
\begin{soalbox}{}
Dalam kotak ada $5$ bola merah dan $3$ bola putih. Satu bola diambil. Tentukan
peluang terambil merah dan peluang \emph{bukan} merah.
\end{soalbox}

\jawab
\langkah{1} Total bola $n(S)=5+3=8$.
\langkah{2} $P(\text{merah})=\dfrac{5}{8}$.
\langkah{3} Komplemen: $P(\text{bukan merah})=1-\dfrac58=\dfrac38$ (yaitu peluang
putih).
\jadi $P(\text{merah})=\dfrac58$, $P(\text{bukan merah})=\dfrac38$.

\begin{jebakan}
\begin{itemize}[leftmargin=*,topsep=2pt]
  \item \textbf{Lupa menjumlah total} ($8$, bukan $5$).
  \item \textbf{$P(\text{bukan merah})=\frac15$?} Gunakan komplemen $1-P$, bukan
  membalik pecahan.
\end{itemize}
\end{jebakan}

\begin{variasi}
Tambah warna ketiga; ambil dua bola (kombinasi, lihat Subbab 8.2).
\end{variasi}

% ---------------- SOAL 8.1.3 ----------------
\soalno{8.1.3}{dadu, kejadian sederhana}
\begin{soalbox}{}
Sebuah dadu dilempar. Berapa peluang muncul mata $>4$?
\end{soalbox}

\jawab
\langkah{1} $S=\{1,2,3,4,5,6\}$, $n(S)=6$.
\langkah{2} Mata $>4$: $\{5,6\}$, $n(A)=2$.
\langkah{3} $P=\dfrac{2}{6}=\dfrac13$.
\jadi peluangnya $\dfrac13$.

\begin{jebakan}
\begin{itemize}[leftmargin=*,topsep=2pt]
  \item \textbf{Memasukkan $4$.} ``$>4$'' \emph{tidak} termasuk $4$; kalau
  ``$\ge4$'' baru $\{4,5,6\}$.
\end{itemize}
\end{jebakan}

\begin{variasi}
Mata $\le2$, mata ganjil, mata prima ($\{2,3,5\}$).
\end{variasi}

% ---------------- SOAL 8.1.4 ----------------
\soalno{8.1.4}{peluang kelipatan}
\begin{soalbox}{}
Dari kartu bernomor $1$--$20$ diambil satu. Berapa peluang terambil kelipatan
$5$?
\end{soalbox}

\jawab
\langkah{1} $n(S)=20$.
\langkah{2} Kelipatan $5$: $\{5,10,15,20\}$, $n(A)=4$.
\langkah{3} $P=\dfrac{4}{20}=\dfrac15$.
\jadi peluangnya $\dfrac15$.

\begin{jebakan}
\begin{itemize}[leftmargin=*,topsep=2pt]
  \item \textbf{Melewatkan $20$.} Kelipatan $5$ hingga $20$ ada $4$
  ($20\div5=4$).
\end{itemize}
\end{jebakan}

\begin{variasi}
Kelipatan $3$; kelipatan $2$ \emph{atau} $3$ (aturan gabungan, Soal 8.1 lanjut).
\end{variasi}

% ---------------- SOAL 8.1.5 ----------------
\soalno{8.1.5}{komplemen langsung}
\begin{soalbox}{}
Peluang hujan $0{,}3$. Berapa peluang tidak hujan?
\end{soalbox}

\jawab
\langkah{1} Komplemen: $P(\text{tidak hujan})=1-P(\text{hujan})=1-0{,}3=0{,}7$.
\jadi $0{,}7$.

\begin{jebakan}
\begin{itemize}[leftmargin=*,topsep=2pt]
  \item \textbf{Total peluang harus $1$.} $P(\text{hujan})+P(\text{tidak})=1$;
  jangan sampai jumlahnya bukan $1$.
\end{itemize}
\end{jebakan}

\begin{variasi}
Tiga kemungkinan (hujan/mendung/cerah) yang jumlahnya $1$; komplemen ``minimal
satu''.
\end{variasi}

% ---------------- SOAL 8.1.6 ----------------
\soalno{8.1.6}{dua dadu, jumlah tertentu}
\begin{soalbox}{}
Dua dadu dilempar. Berapa peluang jumlah mata $=7$?
\end{soalbox}

\begin{ingat}{Ruang sampel dua dadu}
Dua dadu menghasilkan $6\times6=36$ pasangan berurutan (aturan perkalian).
Daftarkan pasangan yang memenuhi syarat.
\end{ingat}

\jawab
\langkah{1} $n(S)=6\times6=36$.
\langkah{2} Jumlah $=7$: $(1,6),(2,5),(3,4),(4,3),(5,2),(6,1)$ --- ada $6$
pasangan.
\langkah{3} $P=\dfrac{6}{36}=\dfrac16$.
\jadi peluangnya $\dfrac16$ (jumlah $7$ adalah yang \emph{paling mungkin} pada
dua dadu).

\begin{jebakan}
\begin{itemize}[leftmargin=*,topsep=2pt]
  \item \textbf{Menghitung $(3,4)$ dan $(4,3)$ sebagai satu.} Dua dadu berbeda
  $\Rightarrow$ pasangan berurutan, keduanya dihitung.
  \item \textbf{$n(S)=11$} (menghitung jumlah $2$--$12$). Ruang sampel adalah
  $36$ \emph{pasangan}, bukan $11$ nilai jumlah.
\end{itemize}
\end{jebakan}

\begin{variasi}
Jumlah $=8$ ($5$ pasangan), jumlah $\ge10$, selisih mata $=1$.
\end{variasi}

% ---------------- SOAL 8.1.7 ----------------
\soalno{8.1.7}{peluang huruf}
\begin{soalbox}{}
Sebuah huruf dipilih acak dari ``MATEMATIKA''. Berapa peluang huruf ``A''?
\end{soalbox}

\jawab
\langkah{1} Hitung total huruf: M-A-T-E-M-A-T-I-K-A $=10$ huruf, jadi $n(S)=10$.
\langkah{2} Hitung huruf ``A'': muncul pada posisi ke-2, ke-6, dan ke-10 $=3$
kali.
\langkah{3} $P(\text{A})=\dfrac{3}{10}$.
\jadi peluangnya $\dfrac{3}{10}$.

\begin{jebakan}
\begin{itemize}[leftmargin=*,topsep=2pt]
  \item \textbf{Menghitung huruf \emph{berbeda} saja} (M,A,T,E,I,K $=6$).
  Penyebut adalah total \emph{semua} huruf ($10$), termasuk yang berulang.
\end{itemize}
\end{jebakan}

\begin{variasi}
Peluang huruf ``M'' ($\frac{2}{10}$), huruf vokal, huruf ``Z'' ($0$).
\end{variasi}

% ---------------- SOAL 8.1.8 ----------------
\soalno{8.1.8}{kejadian saling lepas}
\begin{soalbox}{}
Kejadian $A$ dan $B$ saling lepas, $P(A)=0{,}4$, $P(B)=0{,}3$. Hitung
$P(A\cup B)$.
\end{soalbox}

\begin{ingat}{Saling lepas = tak beririsan}
Jika $A$ dan $B$ tak bisa terjadi bersamaan ($A\cap B=\varnothing$), maka
$P(A\cap B)=0$, sehingga $P(A\cup B)=P(A)+P(B)$ (tanpa pengurangan).
\end{ingat}

\jawab
\langkah{1} Karena saling lepas, $P(A\cap B)=0$.
\langkah{2} $P(A\cup B)=P(A)+P(B)=0{,}4+0{,}3=0{,}7$.
\jadi $P(A\cup B)=0{,}7$.

\begin{jebakan}
\begin{itemize}[leftmargin=*,topsep=2pt]
  \item \textbf{Mengurangi $P(A\cap B)$ padahal $0$.} Untuk saling lepas tak ada
  irisan. Tetapi jika \emph{tidak} saling lepas, wajib kurangi $P(A\cap B)$.
\end{itemize}
\end{jebakan}

\begin{variasi}
Jika tidak saling lepas ($P(A\cap B)=0{,}1$), maka $P(A\cup B)=0{,}6$. Bandingkan.
\end{variasi}

% ---------------- SOAL 8.1.9 ----------------
\soalno{8.1.9}{dua koin, keduanya sama}
\begin{soalbox}{}
Dua koin dilempar: berapa peluang keduanya sama (AA atau GG)?
\end{soalbox}

\jawab
\langkah{1} $S=\{AA,AG,GA,GG\}$, $n(S)=4$.
\langkah{2} ``Keduanya sama'': $\{AA,GG\}$, $n(A)=2$.
\langkah{3} $P=\dfrac{2}{4}=\dfrac12$.
\jadi peluangnya $\dfrac12$.

\begin{jebakan}
\begin{itemize}[leftmargin=*,topsep=2pt]
  \item \textbf{Menganggap ``keduanya sama'' hanya $AA$.} Termasuk juga $GG$.
\end{itemize}
\end{jebakan}

\begin{variasi}
``Keduanya berbeda'' = komplemen = $\frac12$; tiga koin ``ketiganya sama''
($\frac{2}{8}=\frac14$).
\end{variasi}

% ---------------- SOAL 8.1.10 ----------------
\soalno{8.1.10}{peluang bilangan prima}
\begin{soalbox}{}
Dari bilangan $1$--$10$, berapa peluang terambil bilangan prima?
\end{soalbox}

\jawab
\langkah{1} $n(S)=10$.
\langkah{2} Bilangan prima $1$--$10$: $\{2,3,5,7\}$ ($n=4$). Ingat: $1$
\emph{bukan} prima.
\langkah{3} $P=\dfrac{4}{10}=\dfrac25$.
\jadi peluangnya $\dfrac25$.

\begin{jebakan}
\begin{itemize}[leftmargin=*,topsep=2pt]
  \item \textbf{Menghitung $1$ sebagai prima.} Prima harus punya \emph{tepat
  dua} faktor; $1$ hanya punya satu, jadi bukan prima.
  \item \textbf{Lupa $2$ (prima genap satu-satunya).}
\end{itemize}
\end{jebakan}

\begin{variasi}
Bilangan komposit; prima antara $1$--$20$.
\end{variasi}

% ############################################################
\section*{Subbab 8.2 — Pencacahan dan Peluang Kejadian Majemuk}
\addcontentsline{toc}{section}{Subbab 8.2 — Pencacahan}
% ############################################################

\begin{ingat}{Perkalian, permutasi, kombinasi}
\textbf{Aturan perkalian:} tahap-tahap dikalikan. \textbf{Permutasi} (urutan
\emph{penting}): $P(n,r)=\dfrac{n!}{(n-r)!}$. \textbf{Kombinasi} (urutan
\emph{tak penting}): $C(n,r)=\dfrac{n!}{r!(n-r)!}$. Pertanyaan kunci:
\emph{apakah urutan penting?}
\end{ingat}

% ---------------- SOAL 8.2.1 ----------------
\soalno{8.2.1}{bilangan digit berbeda (permutasi)}
\begin{soalbox}{}
Berapa banyak bilangan $3$ digit berbeda yang dapat dibentuk dari angka
$1,2,3,4$?
\end{soalbox}

\jawab
\langkah{1} Urutan penting (angka $123\neq321$) dan tanpa pengulangan
$\Rightarrow$ permutasi $P(4,3)$.
\langkah{2} $P(4,3)=\dfrac{4!}{(4-3)!}=\dfrac{4!}{1!}=4\cdot3\cdot2=24$.
\jadi ada $24$ bilangan.
(\emph{Cara ``kotak'':} $\underline{4}\times\underline{3}\times\underline{2}=24$
--- digit pertama $4$ pilihan, sisanya berkurang.)

\begin{jebakan}
\begin{itemize}[leftmargin=*,topsep=2pt]
  \item \textbf{Memakai kombinasi.} Bilangan \emph{berurut}, jadi permutasi
  ($123$ dan $321$ berbeda).
  \item \textbf{Membolehkan pengulangan.} ``Digit berbeda'' berarti tanpa
  pengulangan; kalau boleh berulang jawabannya $4^3=64$.
\end{itemize}
\end{jebakan}

\begin{variasi}
Boleh berulang ($4^3=64$); dari $0$--$9$ dengan digit pertama tak boleh $0$.
\end{variasi}

% ---------------- SOAL 8.2.2 ----------------
\soalno{8.2.2}{menghitung C dan P}
\begin{soalbox}{}
Hitung $C(6,2)$ dan $P(6,2)$.
\end{soalbox}

\jawab
\langkah{1} $C(6,2)=\dfrac{6!}{2!\,4!}=\dfrac{6\cdot5}{2\cdot1}=15$.
\langkah{2} $P(6,2)=\dfrac{6!}{4!}=6\cdot5=30$.
\jadi $C(6,2)=15$ dan $P(6,2)=30$.

\begin{tips}
Hubungan: $P(n,r)=C(n,r)\cdot r!$. Di sini $30=15\cdot2!$ --- permutasi $=$
kombinasi dikali banyak urutan tiap pilihan.
\end{tips}

\begin{jebakan}
\begin{itemize}[leftmargin=*,topsep=2pt]
  \item \textbf{Lupa $r!$ pada kombinasi.} $C$ membagi tambahan $r!$ untuk
  ``membuang'' urutan.
\end{itemize}
\end{jebakan}

\begin{variasi}
$C(6,4)=C(6,2)=15$ (sifat simetri $C(n,r)=C(n,n-r)$).
\end{variasi}

% ---------------- SOAL 8.2.3 ----------------
\soalno{8.2.3}{ketua \& wakil (permutasi)}
\begin{soalbox}{}
Dari $8$ orang dipilih ketua dan wakil. Berapa cara?
\end{soalbox}

\jawab
\langkah{1} Jabatan berbeda $\Rightarrow$ urutan penting $\Rightarrow$ permutasi.
\langkah{2} $P(8,2)=8\cdot7=56$ (ketua $8$ pilihan, wakil $7$ sisa).
\jadi ada $56$ cara.

\begin{jebakan}
\begin{itemize}[leftmargin=*,topsep=2pt]
  \item \textbf{Memakai kombinasi ($C(8,2)=28$).} Ketua-A wakil-B berbeda dari
  ketua-B wakil-A, jadi urutan \emph{penting}.
\end{itemize}
\end{jebakan}

\begin{variasi}
Memilih ``$2$ anggota komite'' (tanpa jabatan) $\Rightarrow$ kombinasi $C(8,2)=28$.
Bandingkan.
\end{variasi}

% ---------------- SOAL 8.2.4 ----------------
\soalno{8.2.4}{permutasi dengan huruf berulang}
\begin{soalbox}{}
Berapa cara menyusun huruf ``BUKU''?
\end{soalbox}

\begin{ingat}{Permutasi dengan unsur sama}
Jika ada huruf yang \emph{berulang}, bagi dengan faktorial banyak
pengulangannya: $\dfrac{n!}{n_1!\,n_2!\cdots}$. Ini menghapus susunan yang
kelihatan sama.
\end{ingat}

\jawab
\langkah{1} ``BUKU'' punya $4$ huruf, dengan ``U'' berulang $2$ kali.
\langkah{2} Banyak susunan $=\dfrac{4!}{2!}=\dfrac{24}{2}=12$.
\jadi ada $12$ susunan.

\begin{jebakan}
\begin{itemize}[leftmargin=*,topsep=2pt]
  \item \textbf{Menghitung $4!=24$.} Karena dua ``U'' identik, susunan yang
  tampak sama tak boleh dihitung dua kali $\Rightarrow$ bagi $2!$.
\end{itemize}
\end{jebakan}

\begin{variasi}
``MATEMATIKA'' (banyak huruf berulang): $\frac{10!}{3!\,2!\,2!}$ (A$\times3$,
M$\times2$, T$\times2$).
\end{variasi}

% ---------------- SOAL 8.2.5 ----------------
\soalno{8.2.5}{peluang via kombinasi}
\begin{soalbox}{}
Dari $5$ bola merah dan $4$ bola biru diambil $3$ sekaligus. Berapa peluang
ketiganya biru?
\end{soalbox}

\jawab
\langkah{1} ``Sekaligus'' $\Rightarrow$ urutan tak penting $\Rightarrow$
kombinasi. Total cara: $n(S)=C(9,3)=\dfrac{9\cdot8\cdot7}{3\cdot2\cdot1}=84$.
\langkah{2} Ketiganya biru (dari $4$ biru): $n(A)=C(4,3)=4$.
\langkah{3} $P=\dfrac{4}{84}=\dfrac{1}{21}$.
\jadi peluangnya $\dfrac{1}{21}$.

\begin{jebakan}
\begin{itemize}[leftmargin=*,topsep=2pt]
  \item \textbf{Memakai permutasi} untuk ``diambil sekaligus''. Tak ada urutan;
  gunakan kombinasi di pembilang \emph{dan} penyebut.
\end{itemize}
\end{jebakan}

\begin{variasi}
Peluang $2$ merah $1$ biru: $\dfrac{C(5,2)C(4,1)}{C(9,3)}=\dfrac{10\cdot4}{84}
=\dfrac{40}{84}$.
\end{variasi}

% ---------------- SOAL 8.2.6 ----------------
\soalno{8.2.6}{aturan perkalian (PIN)}
\begin{soalbox}{}
Sebuah PIN $4$ digit (boleh berulang). Berapa banyak kemungkinan?
\end{soalbox}

\jawab
\langkah{1} Tiap digit punya $10$ pilihan ($0$--$9$), dan boleh berulang.
\langkah{2} Aturan perkalian: $10\times10\times10\times10=10^4=10000$.
\jadi ada $10.000$ kemungkinan.

\begin{jebakan}
\begin{itemize}[leftmargin=*,topsep=2pt]
  \item \textbf{Memakai permutasi $P(10,4)$.} Itu untuk digit \emph{berbeda}.
  Karena boleh berulang, pakai $10^4$.
\end{itemize}
\end{jebakan}

\begin{variasi}
Digit berbeda ($P(10,4)=5040$); PIN yang harus diawali angka tak-nol.
\end{variasi}

% ---------------- SOAL 8.2.7 ----------------
\soalno{8.2.7}{diagonal segi banyak}
\begin{soalbox}{}
Berapa banyak diagonal segi lima?
\end{soalbox}

\begin{ingat}{Diagonal = ruas antar titik sudut, dikurangi sisi}
Banyak ruas antar $n$ titik $=C(n,2)$; kurangi $n$ sisi. Rumus:
$\dfrac{n(n-3)}{2}$.
\end{ingat}

\jawab
\langkah{1} $n=5$. Total ruas penghubung titik sudut $=C(5,2)=10$.
\langkah{2} Kurangi $5$ sisi: $10-5=5$. (Atau langsung $\dfrac{5(5-3)}{2}
=\dfrac{10}{2}=5$.)
\jadi segi lima punya $5$ diagonal.

\begin{jebakan}
\begin{itemize}[leftmargin=*,topsep=2pt]
  \item \textbf{Lupa mengurangi sisi.} $C(5,2)=10$ mencakup sisi \emph{dan}
  diagonal; sisi bukan diagonal.
\end{itemize}
\end{jebakan}

\begin{variasi}
Segi enam: $\frac{6\cdot3}{2}=9$ diagonal; segi delapan: $20$.
\end{variasi}

% ---------------- SOAL 8.2.8 ----------------
\soalno{8.2.8}{kombinasi memilih}
\begin{soalbox}{}
Dari $10$ soal dipilih $6$. Berapa cara?
\end{soalbox}

\jawab
\langkah{1} Memilih tanpa urutan $\Rightarrow$ kombinasi $C(10,6)$.
\langkah{2} Gunakan simetri $C(10,6)=C(10,4)$ agar lebih mudah:
\[
C(10,4)=\frac{10\cdot9\cdot8\cdot7}{4\cdot3\cdot2\cdot1}=\frac{5040}{24}=210.
\]
\jadi ada $210$ cara.

\begin{jebakan}
\begin{itemize}[leftmargin=*,topsep=2pt]
  \item \textbf{Memakai permutasi.} Memilih soal tak mementingkan urutan
  pengerjaan $\Rightarrow$ kombinasi.
\end{itemize}
\end{jebakan}

\begin{variasi}
``Pilih $6$ dari $10$ dengan $2$ soal wajib'' $\Rightarrow C(8,4)$ (sisanya dari
$8$).
\end{variasi}

% ---------------- SOAL 8.2.9 ----------------
\soalno{8.2.9}{dua dadu, angka kembar}
\begin{soalbox}{}
Dua dadu dilempar. Berapa peluang muncul angka kembar (dua mata sama)?
\end{soalbox}

\jawab
\langkah{1} $n(S)=36$.
\langkah{2} Angka kembar: $(1,1),(2,2),(3,3),(4,4),(5,5),(6,6)$ --- ada $6$.
\langkah{3} $P=\dfrac{6}{36}=\dfrac16$.
\jadi peluangnya $\dfrac16$.

\begin{jebakan}
\begin{itemize}[leftmargin=*,topsep=2pt]
  \item \textbf{Menghitung $(2,2)$ dua kali.} Kembar hanya satu pasangan per
  nilai; ada $6$ total.
\end{itemize}
\end{jebakan}

\begin{variasi}
Peluang \emph{tidak} kembar ($\frac{30}{36}=\frac56$); kembar genap.
\end{variasi}

% ---------------- SOAL 8.2.10 ----------------
\soalno{8.2.10}{tiga koin, tepat dua Gambar}
\begin{soalbox}{}
Tiga koin dilempar. Berapa peluang muncul tepat dua Gambar?
\end{soalbox}

\jawab
\langkah{1} $n(S)=2^3=8$.
\langkah{2} ``Tepat $2$ Gambar dari $3$'': pilih posisi $2$ Gambar $=C(3,2)=3$
(yaitu $GGA,GAG,AGG$).
\langkah{3} $P=\dfrac{3}{8}$.
\jadi peluangnya $\dfrac38$.

\begin{jebakan}
\begin{itemize}[leftmargin=*,topsep=2pt]
  \item \textbf{``Tepat dua'' vs ``minimal dua''.} Tepat $2$ = $3$ kasus;
  minimal $2$ = $3+1$ ($GGG$) $=4$ kasus.
\end{itemize}
\end{jebakan}

\begin{variasi}
Tepat satu Gambar ($\frac38$); tiga Gambar ($\frac18$); pola binomial $C(3,k)/8$.
\end{variasi}

% ############################################################
\section*{Subbab 8.3 — Peluang Lanjutan (Bersyarat, Independen, Bayes)}
\addcontentsline{toc}{section}{Subbab 8.3 — Peluang Lanjutan}
% ############################################################

\begin{ingat}{Bersyarat, independen, Bayes}
\textbf{Bersyarat:} $P(A\mid B)=\dfrac{P(A\cap B)}{P(B)}$ (ruang sampel
menyempit ke $B$). \textbf{Independen:} $P(A\cap B)=P(A)P(B)$. \textbf{Bayes:}
$P(A\mid B)=\dfrac{P(B\mid A)P(A)}{P(B)}$, dengan
$P(B)=P(B\mid A)P(A)+P(B\mid A^c)P(A^c)$.
\end{ingat}

% ---------------- SOAL 8.3.1 ----------------
\soalno{8.3.1}{peluang bersyarat langsung}
\begin{soalbox}{}
Diketahui $P(A\cap B)=0{,}2$ dan $P(B)=0{,}5$. Tentukan $P(A\mid B)$.
\end{soalbox}

\jawab
\langkah{1} $P(A\mid B)=\dfrac{P(A\cap B)}{P(B)}=\dfrac{0{,}2}{0{,}5}=0{,}4$.
\jadi $P(A\mid B)=0{,}4$.

\begin{jebakan}
\begin{itemize}[leftmargin=*,topsep=2pt]
  \item \textbf{Membagi terbalik} ($\frac{P(B)}{P(A\cap B)}$). Syarat ($B$)
  berada di \emph{penyebut}.
\end{itemize}
\end{jebakan}

\begin{variasi}
Diberi $P(A\mid B)$ dan $P(B)$, cari $P(A\cap B)=P(A\mid B)P(B)$.
\end{variasi}

% ---------------- SOAL 8.3.2 ----------------
\soalno{8.3.2}{bersyarat dari data}
\begin{soalbox}{}
Dari $50$ orang: $30$ berkacamata, $20$ pria, $12$ pria berkacamata. Jika
terpilih seorang pria, berapa peluang ia berkacamata?
\end{soalbox}

\begin{ingat}{Ruang sampel menyempit}
``Diketahui pria'' berarti kita hanya melihat $20$ pria. Peluang dihitung
terhadap $20$ itu, bukan $50$.
\end{ingat}

\jawab
\langkah{1} Syaratnya ``pria'' $\Rightarrow$ ruang sampel menyempit ke $20$ pria.
\langkah{2} Di antaranya, $12$ berkacamata.
\langkah{3} $P(\text{kacamata}\mid\text{pria})=\dfrac{12}{20}=0{,}6$.
\jadi peluangnya $0{,}6$.

\begin{jebakan}
\begin{itemize}[leftmargin=*,topsep=2pt]
  \item \textbf{Membagi $50$.} Setelah tahu ``pria'', penyebutnya $20$ (jumlah
  pria), bukan $50$.
\end{itemize}
\end{jebakan}

\begin{variasi}
$P(\text{pria}\mid\text{kacamata})=\frac{12}{30}=0{,}4$ (arah terbalik ---
perhatikan syaratnya).
\end{variasi}

% ---------------- SOAL 8.3.3 ----------------
\soalno{8.3.3}{bersyarat pada dua dadu}
\begin{soalbox}{}
Dua dadu dilempar. Jika diketahui jumlahnya genap, berapa peluang jumlahnya $8$?
\end{soalbox}

\jawab
\langkah{1} Syarat ``jumlah genap'' menyempitkan ruang sampel. Jumlah genap:
$2,4,6,8,10,12$. Banyak pasangan: $1+3+5+5+3+1=18$.
\langkah{2} Jumlah $=8$: $(2,6),(3,5),(4,4),(5,3),(6,2)$ --- ada $5$.
\langkah{3} $P(\text{jumlah }8\mid\text{genap})=\dfrac{5}{18}$.
\jadi peluangnya $\dfrac{5}{18}$.

\begin{jebakan}
\begin{itemize}[leftmargin=*,topsep=2pt]
  \item \textbf{Membagi $36$.} Karena disyaratkan genap, penyebut adalah $18$
  (pasangan berjumlah genap), bukan $36$.
\end{itemize}
\end{jebakan}

\begin{variasi}
``Jika jumlah $>7$, peluang jumlah $10$?''; ``jika dadu pertama $3$, peluang
jumlah $7$?''.
\end{variasi}

% ---------------- SOAL 8.3.4 ----------------
\soalno{8.3.4}{menilai independensi}
\begin{soalbox}{}
Apakah ``hujan hari ini'' dan ``lulus ujian'' independen? Jelaskan.
\end{soalbox}

\jawab
\langkah{1} Dua kejadian independen jika terjadinya satu \emph{tidak mengubah}
peluang yang lain.
\langkah{2} Cuaca hujan tidak memengaruhi hasil ujian (tak ada kaitan sebab
akibat yang wajar).
\jadi \textbf{ya, independen} --- $P(\text{lulus}\mid\text{hujan})=P(\text{lulus})$.
\emph{(Catatan: jika hujan membuat siswa terlambat/terganggu, keterkaitan bisa
muncul; independensi selalu bergantung konteks.)}

\begin{jebakan}
\begin{itemize}[leftmargin=*,topsep=2pt]
  \item \textbf{Mengira ``kejadian berbeda topik'' pasti independen.} Independensi
  adalah soal peluang, bukan sekadar ``tak berhubungan tema''. Uji dengan
  $P(A\mid B)=P(A)$.
\end{itemize}
\end{jebakan}

\begin{variasi}
Contoh \emph{dependen}: ``mendung'' dan ``hujan'' (mendung menaikkan peluang
hujan).
\end{variasi}

% ---------------- SOAL 8.3.5 ----------------
\soalno{8.3.5}{dependen (tanpa pengembalian)}
\begin{soalbox}{}
Dari $4$ merah, $3$ biru diambil $2$ tanpa pengembalian. Berapa peluang keduanya
biru?
\end{soalbox}

\jawab
\textbf{Cara bertahap (dependen).}
\langkah{1} Biru pertama: $\dfrac{3}{7}$.
\langkah{2} Setelah satu biru diambil, tersisa $2$ biru dari $6$ bola:
$\dfrac{2}{6}$.
\langkah{3} $P(BB)=\dfrac{3}{7}\cdot\dfrac{2}{6}=\dfrac{6}{42}=\dfrac17$.

\textbf{Cara kombinasi.} $\dfrac{C(3,2)}{C(7,2)}=\dfrac{3}{21}=\dfrac17$.
\jadi peluangnya $\dfrac17$.

\begin{jebakan}
\begin{itemize}[leftmargin=*,topsep=2pt]
  \item \textbf{Memakai $\frac37\cdot\frac37$.} Itu untuk \emph{dengan}
  pengembalian. Tanpa pengembalian, tahap kedua berubah menjadi $\frac26$.
\end{itemize}
\end{jebakan}

\begin{variasi}
Keduanya merah ($\frac47\cdot\frac36=\frac27$); satu merah satu biru.
\end{variasi}

% ---------------- SOAL 8.3.6 ----------------
\soalno{8.3.6}{independen (dengan pengembalian)}
\begin{soalbox}{}
Sama seperti Soal 8.3.5 tetapi \emph{dengan} pengembalian. Bandingkan hasilnya.
\end{soalbox}

\jawab
\langkah{1} Karena dikembalikan, komposisi bola tak berubah: tiap tahap
$P(\text{biru})=\dfrac37$.
\langkah{2} $P(BB)=\dfrac37\cdot\dfrac37=\dfrac{9}{49}$.
\langkah{3} Bandingkan: tanpa pengembalian $\dfrac17=\dfrac{7}{49}$; dengan
pengembalian $\dfrac{9}{49}$.
\jadi $\dfrac{9}{49}$, sedikit \emph{lebih besar} daripada tanpa pengembalian.
(Masuk akal: dengan pengembalian, biru selalu punya peluang penuh $\frac37$ di
tahap kedua.)

\begin{jebakan}
\begin{itemize}[leftmargin=*,topsep=2pt]
  \item \textbf{Menyamakan kedua kasus.} Pengembalian membuat kejadian
  \emph{independen}; tanpa pengembalian membuatnya \emph{dependen}.
\end{itemize}
\end{jebakan}

\begin{variasi}
Tiga kali dengan pengembalian ($(\frac37)^3$); bandingkan ketiga skenario.
\end{variasi}

% ---------------- SOAL 8.3.7 ----------------
\soalno{8.3.7}{membuktikan independensi}
\begin{soalbox}{}
Sebuah koin dilempar dua kali. Buktikan kedua lemparan independen dengan
menghitung $P(\text{AA})$.
\end{soalbox}

\jawab
\langkah{1} Peluang tiap lemparan Angka: $P(A)=\dfrac12$.
\langkah{2} Jika independen, seharusnya $P(AA)=P(A)\cdot P(A)=\dfrac12\cdot
\dfrac12=\dfrac14$.
\langkah{3} Cek dari ruang sampel: $S=\{AA,AG,GA,GG\}$, $P(AA)=\dfrac14$.
\langkah{4} Karena $P(AA)=P(A)P(A)$, kedua lemparan \textbf{independen}.
\jadi terbukti independen (koin ``tak punya ingatan'').

\begin{jebakan}
\begin{itemize}[leftmargin=*,topsep=2pt]
  \item \textbf{Gambler's fallacy.} Setelah beberapa Angka, peluang Angka
  berikutnya tetap $\frac12$ --- lemparan tak saling memengaruhi.
\end{itemize}
\end{jebakan}

\begin{variasi}
Buktikan $P(AG)=P(A)P(G)=\frac14$; tiga lemparan independen.
\end{variasi}

% ---------------- SOAL 8.3.8 ----------------
\soalno{8.3.8}{diagram pohon koin}
\begin{soalbox}{}
Gambarkan diagram pohon pelemparan koin $2$ kali dan tuliskan peluang tiap
lintasan.
\end{soalbox}

\jawab
Tiap cabang berpeluang $\frac12$; peluang lintasan = hasil kali cabang.

\begin{center}
\begin{tikzpicture}[
  node/.style={circle,draw=biru,fill=birumuda,font=\small\bfseries,inner sep=2pt,minimum size=7mm},
  panah/.style={-Stealth,black!55,thick}]
\node[draw=hijau,fill=hijaumuda,rounded corners,inner sep=4pt] (R) at (0,0) {Mulai};
\node[node] (A) at (3,1.4) {A};
\node[node] (G) at (3,-1.4) {G};
\draw[panah](R)--(A) node[midway,above,font=\scriptsize]{$\tfrac12$};
\draw[panah](R)--(G) node[midway,below,font=\scriptsize]{$\tfrac12$};
\node[node] (AA) at (6,2.3) {A};
\node[node] (AG) at (6,0.6) {G};
\node[node] (GA) at (6,-0.6) {A};
\node[node] (GG) at (6,-2.3) {G};
\draw[panah](A)--(AA) node[midway,above,font=\scriptsize]{$\tfrac12$};
\draw[panah](A)--(AG) node[midway,below,font=\scriptsize]{$\tfrac12$};
\draw[panah](G)--(GA) node[midway,above,font=\scriptsize]{$\tfrac12$};
\draw[panah](G)--(GG) node[midway,below,font=\scriptsize]{$\tfrac12$};
\node[anchor=west,font=\small] at (6.6,2.3){$AA:\tfrac12\cdot\tfrac12=\tfrac14$};
\node[anchor=west,font=\small] at (6.6,0.6){$AG:\tfrac14$};
\node[anchor=west,font=\small] at (6.6,-0.6){$GA:\tfrac14$};
\node[anchor=west,font=\small] at (6.6,-2.3){$GG:\tfrac14$};
\end{tikzpicture}
\end{center}
\jadi keempat lintasan ($AA,AG,GA,GG$) masing-masing berpeluang $\dfrac14$, dan
jumlahnya $1$.

\begin{jebakan}
\begin{itemize}[leftmargin=*,topsep=2pt]
  \item \textbf{Menjumlahkan cabang, bukan mengalikan.} Peluang \emph{sepanjang
  satu lintasan} dikalikan; penjumlahan hanya untuk menggabungkan lintasan
  berbeda.
\end{itemize}
\end{jebakan}

\begin{variasi}
Pohon pengambilan bola tanpa pengembalian (cabang berubah tiap tahap); tiga
lemparan (delapan lintasan).
\end{variasi}

% ---------------- SOAL 8.3.9 ----------------
\soalno{8.3.9}{Teorema Bayes}
\begin{soalbox}{}
Sebuah pabrik: mesin A ($60\%$ produk, cacat $2\%$) dan mesin B ($40\%$, cacat
$5\%$). Sebuah produk cacat; berapa peluang ia dari mesin B?
\end{soalbox}

\begin{ingat}{Bayes = membalik arah}
Kita tahu $P(\text{cacat}\mid B)$; ingin $P(B\mid\text{cacat})$. Hitung dulu
peluang total cacat, lalu bandingkan sumbangan B terhadap total itu.
\end{ingat}

\jawab
\langkah{1} Data: $P(A)=0{,}6,\ P(B)=0{,}4$; $P(\text{cacat}\mid A)=0{,}02$,
$P(\text{cacat}\mid B)=0{,}05$.
\langkah{2} Peluang total cacat:
\[
P(\text{cacat})=0{,}02(0{,}6)+0{,}05(0{,}4)=0{,}012+0{,}020=0{,}032.
\]
\langkah{3} Bayes:
\[
P(B\mid\text{cacat})=\frac{P(\text{cacat}\mid B)P(B)}{P(\text{cacat})}
=\frac{0{,}020}{0{,}032}=0{,}625.
\]
\jadi peluangnya $0{,}625=62{,}5\%$. (Walau B hanya $40\%$ produksi, ia
menyumbang $62{,}5\%$ produk cacat karena tingkat cacatnya lebih tinggi.)

\begin{jebakan}
\begin{itemize}[leftmargin=*,topsep=2pt]
  \item \textbf{Menjawab $0{,}05$ (tingkat cacat B).} Itu $P(\text{cacat}\mid B)$,
  bukan $P(B\mid\text{cacat})$. Bayes membalik arahnya.
  \item \textbf{Lupa peluang total.} Penyebut adalah $P(\text{cacat})$ gabungan
  dari A \emph{dan} B.
\end{itemize}
\end{jebakan}

\begin{variasi}
$P(A\mid\text{cacat})=\frac{0{,}012}{0{,}032}=0{,}375$ (jumlahnya $1$ dengan
B); uji medis (contoh di modul).
\end{variasi}

% ---------------- SOAL 8.3.10 ----------------
\soalno{8.3.10}{base rate \& positif palsu}
\begin{soalbox}{}
Tes narkoba akurat $98\%$, hanya $0{,}5\%$ populasi memakai. Tanpa hitung penuh,
jelaskan mengapa hasil positif belum tentu berarti pemakai.
\end{soalbox}

\begin{ingat}{Masalah base rate}
Ketika kejadian sangat \emph{langka}, positif palsu dari kelompok besar
(non-pemakai) bisa mengalahkan positif benar dari kelompok kecil (pemakai) ---
walau tesnya akurat.
\end{ingat}

\jawab
\langkah{1} Bayangkan $10.000$ orang. Pemakai: $0{,}5\%\times10000=50$ orang;
non-pemakai: $9.950$ orang.
\langkah{2} Positif \emph{benar} (dari $50$ pemakai, akurat $98\%$):
$\approx49$ orang.
\langkah{3} Positif \emph{palsu} (dari $9.950$ non-pemakai, keliru $2\%$):
$\approx199$ orang.
\langkah{4} Dari total $\approx248$ hasil positif, hanya $\approx49$ yang benar
pemakai --- sekitar $20\%$ saja.
\jadi karena \emph{base rate} pemakai sangat kecil, mayoritas hasil positif
justru \textbf{positif palsu}. Hasil positif harus dikonfirmasi tes lanjutan.

\begin{jebakan}
\begin{itemize}[leftmargin=*,topsep=2pt]
  \item \textbf{Mengira ``akurat $98\%$'' berarti $98\%$ positif itu pemakai.}
  Akurasi tes berbeda dari peluang pemakai-jika-positif; base rate langka
  mengubah semuanya (Teorema Bayes).
\end{itemize}
\end{jebakan}

\begin{variasi}
Uji penyakit langka (contoh Bayes di modul); saringan spam. Semua adalah
masalah base rate.
\end{variasi}

% ============================================================
\begin{tips}
\textbf{Ringkasan strategi Bab 8.}
\begin{enumerate}[leftmargin=*,topsep=2pt]
  \item \emph{Tulis ruang sampel} / hitung $n(S)$ dulu. Peluang klasik
  $=\frac{n(A)}{n(S)}$; empiris = frekuensi relatif.
  \item \emph{Komplemen} ($1-P$) ampuh untuk ``minimal''. \emph{Gabungan}
  kurangi $P(A\cap B)$ bila tak saling lepas.
  \item \emph{Pencacahan:} tanya ``urutan penting?'' (permutasi vs kombinasi)
  dan ``boleh berulang?''.
  \item \emph{Lanjutan:} bersyarat menyempitkan ruang sampel; dengan/tanpa
  pengembalian $\to$ independen/dependen; Bayes membalik arah dan mewaspadai
  base rate. \textbf{Hindari gambler's fallacy.}
\end{enumerate}
\end{tips}

\end{document}
